% ===== 原文 pp.132–136（PDF p.142–146） =====
\chapter{3d 拓扑场论}

% [原文 p.132]
\section{3d A-模型}

我们描述一个 3d 拓扑 $\sigma$-模型，其形式与 §3.1 中讨论的 4d A-模型 TFT 类似。这个 3d A-模型 TFT 最早在 [6] 中研究；它可以看作 2d A-模型的 3d 类比。

设 $X$ 是带黎曼度量 $g$ 的实流形，$W$ 是一个 3d 底流形。玻色场由一个映射
\[
  \Phi:W\to X
\]
以及 1-形式
\[
  b\in\Omega^1(W,\Phi^*TX),
\]
和辅助场
\begin{gather*}
  P\in\Omega^0(W,\Phi^*TX),\\
  \tilde P\in\Omega^1(W,\Phi^*TX)
\end{gather*}
组成。费米子由 0-形式
\[
  \beta,\eta\in\Omega^0(W,\Phi^*TX)
\]
以及 1-形式
\[
  \psi,\chi\in\Omega^1(W,\Phi^*TX)
\]
组成。

% [Tutor Note: p.132 把本页场含量与 4d A-模型（p.22）并排对照即可看出这是“维度约化”的痕迹：4d 中成对出现的是 $(\phi,\zeta)$（0-形式）、$(B,\tilde\psi)$（$\Omega^{2+}$）、$(\tilde H,\tilde\chi)$（$\Omega^1$）；3d 中则是 $(\phi,\eta)$（0-形式）、$(b,\psi)$（$\Omega^1$）、$(\beta,P)$（0-形式）、$(\chi,\tilde P)$（$\Omega^1$）。形式次数整体降 1，正是 3d 与 4d 中“自对偶 2-形式 $\leftrightarrow$ 1-形式”的对应。]

% [原文 p.133]
BRST 变分为
\begin{align}
  \delta_Q\phi^I &= \eta^I \notag\\
  \delta_Q\eta^I &= 0 \notag\\
  \delta_Qb^I &= \psi^I-\Gamma^I_{JK}\eta^Jb^K \notag\\
  \delta_Q\psi^I &= \frac12R^I{}_{JKL}b^J\eta^K\eta^L-\Gamma^I_{JK}\eta^J\psi^K \tag{B.1}\\
  \delta_Q\beta^I &= P^I-\Gamma^I_{JK}\eta^J\beta^K \notag\\
  \delta_QP^I &= \frac12R^I{}_{JKL}\beta^J\eta^K\eta^L-\Gamma^I_{JK}\eta^JP^K \notag\\
  \delta_Q\chi^I &= \tilde P^I-\Gamma^I_{JK}\eta^J\chi^K \notag\\
  \delta_Q\tilde P^I &= \frac12R^I{}_{JKL}\chi^J\eta^K\eta^L-\Gamma^I_{JK}\eta^J\tilde P^K \notag
\end{align}
其中 $\Gamma^I_{JK}$ 是 Christoffel 符号，$R^I{}_{JKL}$ 是关于度量 $g_{IJ}$ 的 Levi-Civita 联络的 Riemann 张量分量。BRST-正合作用量为
\[
  S=\frac{1}{e^2}\delta_Q\int_W\dd^3x\sqrt{h}\,g_{IJ}\left\{\chi^{Ir}\left(\tilde P_r{}^J-2D_r\phi^J+2\sqrt{h}\,\varepsilon_{rst}D^sb^{Jt}\right)+\beta^I\left(P^J-2D_rb^{Jr}\right)\right\},
\]
其中 $r,s,t=1,2,3$ 是 $W$ 上的形式指标，$D$ 是 $X$ 上 Levi-Civita 联络拉回到 $W$ 的拉回。如 [6] 所述，局部可观测量对应靶空间 $X$ 的上同调类，线可观测量对应 $X$ 上平坦联络的和乐在 $\sigma$-模型映射下的拉回。此外，最简单的一类“几何”边界条件对应闭子流形 $Y\subset X$。注意，由于我们没有假设 $X$ 带有任何额外的几何结构（例如辛形式或复结构），$Y$ 在 $X$ 中的嵌入没有额外条件。

% [Tutor Note: p.133 作用量中 $2\sqrt{h}\,\varepsilon_{rst}D^sb^{Jt}$ 的 $\sqrt h$ 因子可用 3d 的 Hodge 星核对：对 1-形式 $b$ 有 $(\star\dd b)_r=\frac12\sqrt{h}\,\varepsilon_{rst}(\dd b)^{st}=\sqrt{h}\,\varepsilon_{rst}\partial^sb^t$，故 $\sqrt h$ 的幂次正确（若误写成 $1/\sqrt h$ 则与 $g_{IJ}$ 缩并后的量纲不符）。]

\section{带规范场的 3d A-模型}

现在假设靶空间 $X$ 容许李群 $G$ 通过等距作用。我们希望把以 $X$ 为靶空间的 3d A-模型与一个 3d 拓扑规范理论耦合，从而对它做规范化，做法与 §3.2 中对 4d A-模型所做的大致相同。事实上，我们将把 3d A-模型与一个我们称之为 A-型 3d 规范理论的理论耦合。

\subsection{A-型 3d 拓扑规范理论}

% [原文 p.134]
在 3d 中有两种不同的拓扑规范理论，可以通过扭转 $\cN=4$、$d=3$ SYM 理论得到。第一种是 $\cN=2$、$d=4$ SYM 理论的 Donaldson-Witten 扭转的维数约化（见 §3.2.1）。第二种是 3d 内禀的，最早由 Blau 与 Thompson [29] 讨论。我们分别称它们为 A-型与 B-型拓扑规范理论。（采用这一术语的原因是：前者的 BPS 方程是椭圆的，与通常的 A-模型一样；而后者的 BPS 方程是超定的，与通常的 B-模型一样。）

A-型 3d 规范理论的玻色场含量由一个 3-流形 $W$ 上主 $G$-丛 $P$ 上的联络 $A$ 以及 0-形式
\begin{gather*}
  \sigma\in\Omega^0(W,\mathit{ad}\,P)\otimes\C\\
  \varsigma\in\Omega^0(W,\mathit{ad}\,P)
\end{gather*}
和辅助 1-形式
\[
  H\in\Omega^1(W,\mathit{ad}\,P)
\]
组成。费米子场含量由形式
\begin{gather*}
  \rho,\tilde\rho\in\Omega^0(W,\mathit{ad}\,P)\\
  \lambda,\tilde\lambda\in\Omega^1(W,\mathit{ad}\,P)
\end{gather*}
组成。BRST 变换由下式给出 [6]：
\begin{gather*}
  \delta_QA=\lambda,\qquad \delta_Q\lambda=-\dd_A\sigma,\\
  \delta_Q\sigma=0,\\
  \delta_Q\bar\sigma=\tilde\rho,\qquad \delta_Q\tilde\rho=[\sigma,\bar\sigma],\\
  \delta_Q\tilde\lambda=H,\qquad \delta_QH=[\sigma,\tilde\lambda],\\
  \delta_Q\varsigma=\rho,\qquad \delta_Q\rho=[\sigma,\varsigma].
\end{gather*}
作用量为
\[
  S=\frac{1}{e^2}\delta_Q\int_W\dd^3x\sqrt{h}\,\Tr\left\{4\tilde\lambda^r\left(H_r-D_r\varsigma-(\star F)_r\right)+2\bar\sigma D_r\lambda^r\right\},
\]
其中 $F$ 是 $A$ 的曲率，$\star$ 是 3d Hodge 星。

% [Tutor Note: p.134 两种 3d 拓扑规范理论的根本差别可在作用量层面核对：A-型的作用量（上式）整体是 BRST-正合的；而 B-型的 (B.4) 却含一个非 BRST-正合的“拓扑项” $\int_W\Tr(\rho\wedge\dd_{\mathcal A}\zeta)$（见 p.137）。这也解释了为什么 §B.4 要讨论 $\rho$ 取值于 $\fg^*$ 而非 $\fg$（拓扑项要写成与度量无关的形式）。]

\subsection{对 3d A-模型做规范化}

设 $V_a$（$a=1,2,\cdots,\dim G$）是 $X$ 上对应 $G$-作用生成元的向量场。物质扇区的变分如何修改，与 §3.2 中对带规范场 4d A-模型所做的十分相似：

% [原文 p.135]
\begin{align}
  \delta_Q\phi^I &= \eta^I \notag\\
  \delta_Q\eta^I &= \sigma^aV^I_a \notag\\
  \delta_Qb^I &= \psi^I-\Gamma^I_{JK}\eta^Jb^K \notag\\
  \delta_Q\psi^I &= \frac12R^I{}_{JKL}b^J\eta^K\eta^L-\Gamma^I_{JK}\eta^J\psi^K+\sigma^ab^J\nabla_JV^I_a \tag{B.2}\\
  \delta_Q\beta^I &= P^I-\Gamma^I_{JK}\eta^J\beta^K \notag\\
  \delta_QP^I &= \frac12R^I{}_{JKL}\beta^J\eta^K\eta^L-\Gamma^I_{JK}\eta^JP^K+\sigma^a\beta^J\nabla_JV^I_a \notag\\
  \delta_Q\chi^I &= \tilde P^I-\Gamma^I_{JK}\eta^J\chi^K \notag\\
  \delta_Q\tilde P^I &= \frac12R^I{}_{JKL}\chi^J\eta^K\eta^L-\Gamma^I_{JK}\eta^J\tilde P^K+\sigma^a\chi^J\nabla_JV^I_a \notag
\end{align}
其中
\[
  \nabla_JV^I_a=\partial_JV^I_a+\Gamma^I_{JK}V^K_a.
\]
作用量由把物质扇区作用量与规范扇区作用量相加、并把导数规范协变化而得到：
\[
\begin{aligned}
  S=\frac{1}{e^2}\delta_Q\int_W\dd^3x\sqrt{h}\,&\Tr\left\{4\tilde\lambda^r\left(H_r-D_r\varsigma-(\star F)_r\right)+2\bar\sigma D_r\lambda^r\right.\\
  &\left.+\chi^{Ir}\left(\tilde P_r{}^J-2\mathcal{D}_r\phi^J+2\sqrt{h}\,\varepsilon_{rst}\mathcal{D}^sb^{Jt}\right)+\beta^I\left(P^J-2\mathcal{D}_rb^{Jr}\right)\right\}
\end{aligned}
\]
其中 $\mathcal{D}$ 是规范协变导数，与第 3 章中对 4d A-模型讨论的相同。我们指出，这个带规范场 3d A-模型的作用量远非唯一，因为可以加入不影响理论的非极小 BRST-正合项。此外，3d A-模型除与 3d A-型规范理论耦合外，还可以与其它规范扇区耦合来规范化；例如在 §5.2 中我们写下这样一个理论，其规范扇区是把李代数取值的标量 $\varsigma$ 换成一个李群取值的标量而得到的。

% [Tutor Note: p.135 (B.2) 与 4d 的 (3.6) 逐行对照即可确认规范化模版完全相同：首行 $\delta_Q\eta^I=\sigma^aV^I_a$ 对应 (3.6) 的 $\delta_Q\zeta^I=\sigma^aV^I_a$；末项均以 $\sigma^a(\cdot)\nabla_JV^I_a$ 的形式出现。核对要点是 $\nabla$ 必须用协变导数（含 $\Gamma^I_{JK}V^K_a$），否则 $\delta_Q^2=\delta_\sigma$ 不成立。]

% [Tutor Note: p.135 原文“replacing the Lie algebra-valued scalar $\varsigma$ by a a Lie group-valued scalar”有重复的 “a”（笔误，无实质影响）；译文按“一个李群取值的标量”处理。]

\section{Rozansky-Witten 模型}

我们描述一个 3d 拓扑 $\sigma$-模型，它是 2d B-模型的 3d 类比。Rozansky-Witten $\sigma$-模型 [14] 对任意超 Kähler 靶空间 $X$ 与 3d 底流形 $W$ 都有定义。

其玻色场由一个映射
\[
  \Phi:W\to X
\]
组成，我们用 $X$ 上的局部复坐标 $\phi^i$ 与 $\phi^{\bar\imath}$ 来描述它。

% [原文 p.136]
其费米子场由 0-形式
\[
  \eta\in\Omega^0(W,\Phi^*T^{0,1}X)
\]
以及 1-形式
\[
  \chi\in\Omega^1(W,\Phi^*T^{1,0}X)
\]
组成。BRST 变分为
\begin{gather*}
  \delta_Q\phi^i=0\\
  \delta_Q\phi^{\bar\imath}=\eta^{\bar\imath}\\
  \delta_Q\eta^{\bar\imath}=0\\
  \delta_Q\chi^i=\dd\phi^i.
\end{gather*}
BRST-正合作用量为
\[
  S=\delta_Q\int_W\left(g_{i\bar\jmath}\chi^i\wedge\star\dd\phi^{\bar\jmath}
  +\frac12\Omega_{ij}\chi^i\wedge D\chi^j
  +\frac16\Omega_{ij}R^j{}_{kl\bar m}\chi^i\wedge\chi^j\wedge\chi^k\wedge\chi^i\wedge\eta^{\bar m}\right),
\]
其中 $D\chi^i=\dd\chi^i+\Gamma^i_{jk}\dd\phi^j\wedge\chi^k$ 是 $X$ 上 Levi-Civita 联络的拉回，$\star$ 是 3d Hodge 星。

% [Tutor Note: p.136 由 BRST 变分可见 $\eta^{\bar\imath}=\delta_Q\phi^{\bar\imath}$ 是 BRST-正合的，而 $\delta_Q\eta^{\bar\imath}=0$：故 $\eta$ 在局部可观测量的上同调中不出现，局部可观测量只能由全纯坐标 $\phi^i$ 构造——这与原文“局部可观测量对应 Dolbeault 上同调类”的断言一致，可作为校验。]

% [Tutor Note: p.136 为什么靶空间必须“超 Kähler”而不是任意复流形？核对作用量用到的几何数据即可：$g_{i\bar\jmath}$（Kähler 度量）与全纯辛形式 $\Omega_{ij}$（用于 $\frac12\Omega_{ij}\chi^i\wedge D\chi^j$ 与曲率项）以及复结构（用于 $T^{1,0}$、$T^{0,1}$ 的分解）。超 Kähler 条件保证 $\Omega$ 全纯闭且与 Kähler 度量相容；若 $X$ 只是复流形而无全纯辛形式，$\Omega_{ij}$ 与曲率项就写不出来。]

% [Tutor Note: p.136 作用量最后一项的 $\chi$ 指标有排印错误：印作 $\chi^i\wedge\chi^j\wedge\chi^k\wedge\chi^i$（共 4 个 $\chi$），而 $W$ 是 3-流形，4 个 1-形式的外积恒为零；且 $\Omega_{ij}$ 关于 $i,j$ 反对称使 $\Omega_{ij}\chi^i\wedge\chi^j=0$。按形式次数与指标缩并判断，该项应为 $\Omega_{ij}R^j{}_{kl\bar m}\chi^i\wedge\chi^k\wedge\chi^l\wedge\eta^{\bar m}$（3 个 $\chi$ + 1 个 $\eta$，恰为 3-形式）。译文照原文排印，此处指出。]

如 [36] 所述，Rozansky-Witten 模型的局部可观测量对应 $X$ 的 Dolbeault 上同调类。

Rozansky-Witten 模型最简单的边界条件 [36] 对应靶空间 $X$ 的复、拉格朗日子流形。此外，人们可以把 Rozansky-Witten 模型与定义在 2d 边界上的 B-模型耦合；这样的边界条件对应复、拉格朗日子流形上的 Calabi-Yau 纤维化。细节见 [36]。

\section{B-型 3d 拓扑规范理论}

我们描述 B-型 3d 规范理论，它最早由 Blau 与 Thompson [29] 构造。其玻色场含量由一个 3-流形 $W$ 上主 $G$-丛 $P$ 上的联络 $A$ 以及 1-形式
\[
  \varphi\in\Omega^1(W,\mathit{ad}\,P)
\]
和辅助场
\[
  P\in\Omega^0(W,\mathit{ad}\,P)
\]
组成。
